\section*{Глава 1. Анализ предметной области}
\addcontentsline{toc}{section}{Глава 1. Анализ предметной области}
\stepcounter{section}

\subsection{Обзор существующих платформ}

В области платформ для подготовки к олимпиадам по программированию можно выделить следующие решения~\cite{petrov2024}:~\ref{tab:competitors}.

\begin{table}[ht!]
	\centering
	\caption{Сравнительный анализ популярных платформ}
	\label{tab:competitors}
	
	\small
	\resizebox{\textwidth}{!}{
	\begin{tabular}{|l|p{0.3\textwidth}|c|c|c|c|}
		\hline
		\textbf{Название} & \textbf{Краткое описание} & \textbf{Русский язык} & \textbf{Индивидуальные планы} & \textbf{ИИ ассистент} & \textbf{Удобный интерфейс} \\ \hline
		
		LeetCode & Платформа для подготовки к техническим собеседованиям с большим количеством задач & Нет & Нет & Частично & Да \\ \hline
		
		CodeRun & Образовательная платформа с разбором задач и рекомендациями на основе ИИ & Да & Нет & Да & Да \\ \hline
		
		CodeForces & Платформа для спортивного программирования и участия в соревнованиях & Частично & Нет & Нет & Нет \\ \hline
		
		Информатикс & Российская образовательная платформа с задачами по программированию разного уровня & Да & Нет & Нет & Нет \\ \hline
		
		CodeWars & Платформа с геймификацией обучения через задачи (kata) & Нет & Нет & Нет & Да \\ \hline
		
	\end{tabular}
	}
\end{table}

Анализ показал, что существующие платформы обладают рядом преимуществ, однако отсутствует единое решение, полностью удовлетворяющее потребности русскоязычных пользователей.

\subsection{Анализ программных инструментов разработки веб-приложений}

Для реализации проекта были рассмотрены различные современные подходы и технологии разработки веб-приложений.

\begin{itemize}
	\item \textbf{Frontend-фреймворк:} В качестве технологии для разработки клиентской части рассматривались такие решения, как \textit{React}, \textit{Vue} и \textit{Angular}. Angular предоставляет полноценную архитектуру "из коробки", однако является более сложным и избыточным для данного проекта. Vue отличается простотой освоения, но уступает React по распространённости и количеству доступных библиотек. В результате был выбран React~\cite{react_docs} в сочетании с Vite. Данный стек обеспечивает высокую производительность за счёт виртуального DOM, гибкость архитектуры благодаря компонентному подходу, а также быструю разработку благодаря мгновенной сборке и горячей перезагрузке, реализуемой Vite.

	\item \textbf{Backend-платформа:} Среди серверных технологий были рассмотрены Django, Flask и Node.js~\cite{cantelon2017node} (Express/NestJS). Flask и Express предоставляют большую гибкость, но требуют значительной самостоятельной реализации базовой функциональности (аутентификация, администрирование, структура проекта). Node.js хорошо подходит для высоконагруженных асинхронных систем, однако требует дополнительной настройки и архитектурной дисциплины. Django был выбран благодаря наличию встроенных компонентов (ORM, система аутентификации, административная панель), высокой степени безопасности и чёткой архитектуре, что существенно ускоряет разработку и снижает вероятность ошибок.

    \item \textbf{Система управления базой данных:} Рассматривались PostgreSQL, MySQL и SQLite. SQLite подходит для небольших проектов и прототипирования, но не обеспечивает необходимой масштабируемости. MySQL обладает высокой производительностью, однако уступает PostgreSQL в поддержке сложных запросов и расширенных возможностей (например, работа с JSON, транзакции, аналитические функции). PostgreSQL была выбрана как наиболее надёжная и функционально богатая СУБД, обеспечивающая эффективную работу с большими объёмами данных и хорошую интеграцию с Django.

    \item \textbf{Контейнеризация:} В качестве альтернатив Docker рассматривались виртуальные машины и ручная настройка окружения. Использование виртуальных машин требует больше ресурсов и времени на настройку, а ручная конфигурация увеличивает риск несовместимости окружений. Docker позволяет изолировать компоненты системы, стандартизировать окружение и упростить процесс развертывания, что особенно важно при командной разработке и переносе проекта между различными средами.

	\item \textbf{Искусственный интеллект:} Для реализации интеллектуального ассистента рассматривались различные API (например, OpenAI, локальные модели и DeepSeek). Локальные модели требуют значительных вычислительных ресурсов и сложной настройки. OpenAI предоставляет высокое качество, но может быть менее выгоден с точки зрения стоимости. DeepSeek API был выбран как компромисс между качеством, стоимостью и удобством интеграции, обеспечивая достаточный уровень функциональности для задач проекта.
\end{itemize}

\subsection{Формулировка цели и задач работы}

На основе проведенного анализа и требований, изложенных в нормативных документах~\cite{gost_report}, сформулирована цель исследования.

\textbf{Цель работы} --- создание веб-приложения для подготовки к олимпиадам по программированию.

Для достижения поставленной цели необходимо решить следующие \textbf{задачи}:
\begin{enumerate}
	\item Провести анализ предметной области;
	\item Спроектировать архитектуру системы;
	\item Разработать структуру базы данных;
	\item Реализовать интеллектуальный анализ решений;
	\item Обеспечить хранение данных;
	\item Провести тестирование системы.
\end{enumerate}